% ---------------------------------------------------------
%  PROJECT DOCUMENTATION (XeLaTeX SAFE VERSION)
% ---------------------------------------------------------

\documentclass[12pt,a4paper]{article}

% ---------------------------------------------------------
%  LANGUAGE (polyglossia must go BEFORE fontspec!)
% ---------------------------------------------------------
\usepackage{polyglossia}
\setmainlanguage{russian}
\setotherlanguage{english}

% ---------------------------------------------------------
%  FONTS (DejaVu has full Unicode support)
% ---------------------------------------------------------
\usepackage{fontspec}
\setmainfont{DejaVu Serif}
\setsansfont{DejaVu Sans}
\setmonofont{DejaVu Sans Mono}

% Unicode fallback for symbols (─ │ └ ← → etc.)
\usepackage{ucharclasses}
\setTransitionsForSymbols{\fontspec{DejaVu Sans Mono}}{}

\defaultfontfeatures{Ligatures=TeX}

% ---------------------------------------------------------
%  PACKAGES
% ---------------------------------------------------------
\usepackage{geometry}
\usepackage{hyperref}
\usepackage{titlesec}
\usepackage{graphicx}
\usepackage{setspace}
\usepackage{listings}
\usepackage{xcolor}

\geometry{margin=2.5cm}
\onehalfspacing

% ---------------------------------------------------------
%  CODE BLOCK STYLE
% ---------------------------------------------------------
\lstset{
    basicstyle=\ttfamily\small,
    breaklines=true,
    frame=single,
    keywordstyle=\color{blue},
    commentstyle=\color{gray},
    showstringspaces=false,
    tabsize=2
}

% ---------------------------------------------------------
%  DOCUMENT START
% ---------------------------------------------------------
\begin{document}

% -------------------------------
%   TITLE PAGE
% -------------------------------
\begin{titlepage}
    \centering
    {\Large \textbf{Tetris Game Project}}\\[10pt]
    {\large Course: C Core + Makefiles + GCOV + Unit Testing}\\[10pt]
    {\large Author: \textbf{Ivelissc}}\\[10pt]
    {\large Date: \today}\\[40pt]
    \vfill
\end{titlepage}

% -------------------------------
%   TABLE OF CONTENTS
% -------------------------------
\tableofcontents
\newpage

% ---------------------------------------------------------
\section{Overview}

Данный документ описывает структуру, архитектуру и реализацию \textbf{проекта игры Тетрис}.
Он включает в себя информацию о:

\begin{itemize}
    \item backend logic (game engine)
    \item frontend (CLI interface)
    \item Makefile structure
    \item test suite
    \item gcov/lcov coverage report
    \item installation and distribution
\end{itemize}

Проект следует стилю кодирования S21 и использует модульную архитектуру.

% ---------------------------------------------------------
\section{Project Structure}

\begin{lstlisting}
project/
|
+-- brick_game/
|   +-- tetris/
|       +-- backend.c
|       +-- backend.h
|       +-- figures.c
|       +-- figures.h
|
+-- gui/
|   +-- cli/
|       +-- main.c
|       +-- frontend.c
|       +-- frontend.h
|
+-- tests/
|   +-- test_tetris.c
|
+-- Makefile
\-- dvi/
    \-- documentation.tex  <- this file
\end{lstlisting}

\section{State Machine Diagram}

На рисунке ниже представлена диаграмма конечного автомата игры Тетрис.

\begin{figure}[h!]
    \centering
    \includegraphics[width=\textwidth]{fsm_diagram.png}
    \caption{Диаграмма конечного автомата игры}
\end{figure}

% ---------------------------------------------------------
\section{Backend Architecture}

Backend реализует основную игровую механику:

\begin{itemize}
    \item игровое поле;
    \item текущая и следующая фигуры;
    \item обработку ввода;
    \item перемещение и вращение фигур;
    \item уничтожение заполненных линий;
    \item подсчёт очков и уровней;
    \item завершение игры.
\end{itemize}

\subsection{Game State}

\begin{lstlisting}[language=C]
typedef struct {
    int field[FIELD_HEIGHT][FIELD_WIDTH];
    int current[NEXT_H][NEXT_W];
    int next[NEXT_H][NEXT_W];
    int score;
    int high_score;
    int level;
    int speed;
    bool paused;
    long long last_tick_ms;
    int cur_row;
    int cur_col;
    GameEngineState state;
} InternalState;
\end{lstlisting}

% ---------------------------------------------------------
\section{Frontend (CLI)}

Отрисовка выполняется с помощью **ncurses** и включает:

\begin{itemize}
    \item игровое поле;
    \item следующую фигуру;
    \item счёт и уровень;
    \item состояние паузы и конца игры.
\end{itemize}

% ---------------------------------------------------------
\section{Makefile Structure}

Цели Makefile:

\begin{itemize}
    \item \verb|all| - собрать игру
    \item \verb|install| - установить бинарник и библиотеку
    \item \verb|uninstall| - удалить из install dir
    \item \verb|test| - запустить юнит-тесты
    \item \verb|gcov_report| - получить отчёт покрытия
    \item \verb|clean| - очистить build/
    \item \verb|dist| - создать архив проекта
    \item \verb|dvi| - собрать PDF-документацию
\end{itemize}

Пример цели dvi:

\begin{lstlisting}[language=make]
dvi:
    mkdir -p dvi
    xelatex -output-directory=dvi dvi/documentation.tex
\end{lstlisting}

% ---------------------------------------------------------
\section{Unit Testing}

Тестирование реализовано через библиотеку **Check**.
Покрытие должно составлять не менее 80\%.

\subsection{Пример теста}

\begin{lstlisting}[language=C]
START_TEST(init_1) {
    initGame();
    GameInfo_t gi = updateCurrentState();
    ck_assert_int_eq(gi.score, 0);
    closeGame();
}
END_TEST
\end{lstlisting}

% ---------------------------------------------------------
\section{Game Mechanics}

\begin{itemize}
    \item вращение фигур;
    \item перемещение по горизонтали;
    \item ускорение падения (soft drop);
    \item мгновенное падение (hard drop);
    \item уничтожение заполненных линий;
    \item спавн новой фигуры;
    \item завершение игры при достижении верха поля.
\end{itemize}

\subsection{Управление}

\begin{itemize}
    \item left - движение влево
    \item right - движение вправо
    \item up - ускорение падения
    \item Пробел - вращение
    \item P - пауза
    \item Q - выход
\end{itemize}

% ---------------------------------------------------------
\end{document}
